\scanpage{113}
\uline{Proposition} 4.2 \deleted{\unsure{a)}} a) Stabilité de \textquotedblleft eff$^{\mathrm t}$. adm. de poids $\subset P$ et niveau $\leq\unsure{\nu|v}$\textquotedblright\ par sous-faisceau constr., quotients constructibles, sommes finies, extensions.

b) $F_i$ \textquotedblleft eff. adm de poids $\subset P_i$, niveau $\leq\nu_i$\textquotedblright\ pour $i=1,2$ implique $F_1\otimes F_2$ \textquotedblleft eff$^{\mathrm t}$. adm. de poids $\subset P_1+P_2$, niveau $\leq\nu_1+\nu_2$\textquotedblright.

c) Stabilité par chang$^{\mathrm{mt}}$ de base.

d) Si $F$ \inserted{eff$^{\mathrm t}$} admissible \uline{constant tordu} \textquotedblleft à poids $\deleted{\unsure{\leq q|k}}\inserted{\nu+2k}$ et niveau $\leq\nu$\textquotedblright, \deleted{et si}\inserted{(où $k\geq0$, $\nu\geq0$)}, alors pour tout $n\geq\nu+k$\inserted{ (soit $n=\nu+k+h$, $h\geq0$)}, $\check F(-n)$ est \textquotedblleft eff$^{\mathrm t}$. admissible constant tordu, à poids $2n-\nu-2k=\nu+2h$, niveau $\leq\nu$\textquotedblright.

\uline{\unsure{Remarks|Remarq|Remarque}}: introduire plutôt: \textquotedblleft à morceaux de poids $i_\alpha$, niveau $\nu_\alpha$\textquotedblright, $\alpha\in\ldots$\textquotedblright.

\margin{\uline{Utiliser Lefschetz} \uline{ou Deligne}}
% slanted
\uline{Théorème} 5.1. Soit $f:X\longrightarrow S$ morphisme séparé de \deleted{\unsure{t.f.|type}}\inserted{prés finie}, et $F$ \deleted{eff$^{\mathrm t}$ admissible \unsure{de niveau de poids} \unsure{facteur}}\inserted{\uline{strictement spécial} de poids $j$} sur $X$ (i.e. déf. 1 s'applique, avec $F$ \uline{facteur direct} d'un $R^jg_*(Q_{lY})$, où $g:Y\longrightarrow X$ \deleted{\unsure{de}}\inserted{proj. lisse}). \deleted{\unsure{(i.e. déf. 2.1, est de poids $\subset P$, niveau|(i.e. déf. 2.1, avec} \illegible{2}} Alors $R^if_!(F)$ est \uline{eff$^{\mathrm t}$ admissible} \deleted{\unsure{constructible}}\inserted{à poids $\leq$\deleted{\unsure{q}}$i+j$ \deleted{\unsure{et}} et niveaux $\leq\operatorname{Inf}(i+j,2n-i+j)$} \deleted{poids $\unsure{\leq i+j}$ \unsure{si $f$ est lisse et $X$ est} \unsure{de dim. relative $n$}, de niveau \unsure{$i+j$}. Si $n=\operatorname{Sup}\dim X_s$, et si $i\geq n$, alors les poids}.

\margin{$+$ résolution des singularités forte sur les schémas parfaits, aux corps résiduels de $S$.}
% slanted
\margin{$+$ $\frac12$ simplicité \unsure{arithmétiques|arithmétique} ?}
% slanted
Si $n=\operatorname{Sup}_{s\in S}\dim X_s$, \deleted{\unsure{ses}} et si $i\geq n$, les poids \deleted{\unsure{et}} et niveaux \deleted{\unsure{sont}} possibles sont
\diagram{\[
\begin{array}{c|c|c|c|c|}
\text{poids}&2(i-n)&2(i-n)+1&\ldots&(i+j)\\ \hline
\text{niveaux}&\leq0&\leq1&\ldots&\leq2n-i+j
\end{array}
\]}

\uline{Démonstration} On \deleted{\unsure{en|(a)}} peut supposer $F=R^jg_*(Q_{lY})$, $g:Y\longrightarrow X$ proj. lisse. \deleted{\illegible{2}} Utilisant le th. de chang$^{\mathrm{mt}}$ de base propre, \deleted{\illegible{2}} et dévissage via suite exacte de cohomologie à support propre, on est ramené au cas où \deleted{$\exists\, S'\longrightarrow S$} $S$ est affine, \unsure{nell$^{\mathrm e}$|nlle|noeth.}, intègre, $\exists\, S'\longrightarrow S$ morphisme fini plat radiciel, \inserted{$X$ affine,} tel que $X'=X\times_S S'$ est \uline{lisse} sur $S'$, et se \textquotedblleft compactifie\textquotedblright\ en un $\widehat X_{S'}$ \deleted{tel que $X'$} projectif et lisse, tel que $\widehat X_{S'}-X_{S'}$ soit un diviseur à croisements normaux relatifs, à \unsure{compos$^{\mathrm{tes}}$|composantes} irréductibles \uline{lisses} sur $S'$; \inserted{enfin $F$ \deleted{\unsure{sur}} \illegible{1}} \inserted{étant \uline{spécial} de poids $j$.} Quitte à changer de notation, on peut supposer $S=S'$, \deleted{et} d'où $X\subset X'=\hat X$. On trouvera alors, \inserted{si $F=Q_{lX}$,} comme dans les notes préliminaires (partie A, 1)). [\uline{NB} On n'utilise pas la dégénérescence de la suite spectrale qui y intervient.]
\margin{à expliciter ce point en utilisant successiv$^{\mathrm{mt}}$ \unsure{corollaire de} 3.3. a) et 4.2 a)}
% slanted
% archive sheet number: 112
